\documentclass[10pt,a4paper]{amsart}
\usepackage[margin=19mm]{geometry}
\usepackage{amsmath,amssymb,mathtools}
\usepackage[scr=rsfs,cal=euler]{mathalfa}
\usepackage{xfrac}
\usepackage[unicode,hidelinks,bookmarks=false]{hyperref}
\newcommand{\ie}{i.e.\ }
\newcommand{\cf}{cf.\ }
\newcommand{\resp}{respectively\ }
\newcommand{\Subsection}{\subsection}
\providecommand{\nobreakdash}{\nobreak\hbox{-}}
\setlength{\emergencystretch}{2em}
\multlinegap0pt
\usepackage{amssymb}
\usepackage{dsfont}
\usepackage{float}
\usepackage{bm}
\usepackage{mathtools}
\usepackage[shortlabels]{enumitem}
\usepackage{tikz}
\newtheorem{lemma}{Lemma}
\renewcommand{\H}{\mathcal H}
\newcommand{\R}{\mathbb R}
\title{Low regularity potentials in heterogeneous Cahn--Hilliard functionals}
\author{Riccardo Cristoferi, Jakob Deutsch, Luca Pignatelli}
\date{Source: arXiv:2603.25369v1 (CC BY 4.0)}
\begin{document}
\maketitle

\noindent\emph{Source and licence.} Low regularity potentials in heterogeneous Cahn--Hilliard functionals. Version arXiv:2603.25369v1 (CC BY 4.0), distributed by arXiv under the Creative Commons Attribution 4.0 International licence, http://creativecommons.org/licenses/by/4.0/, as recorded in the arXiv licence metadata for this exact version and shown on the abstract page https://arxiv.org/abs/2603.25369v1. Full licence notice: \texttt{licenses/arxiv-2603.25369-CC-BY-4.0.txt}.\par\smallskip

\noindent\textit{Editorial scope.} Counted unit: the article's lemma lemma:technical\_approx\_limits (source0.tex lines 2581--2591). The article's complete original proof of it is delivered with it (source0.tex lines 2593--2619, one \texttt{proof} environment of 27 lines). No mathematics is written, repaired or re-derived: the statement and the proof are the article's own lines, in the article's own notation, and no retained line is substituted or re-flowed.  No further source-local context is needed: every cross-reference of every retained line resolves inside the two delivered spans.  Nothing was omitted from any retained span, so the package carries no omission record.  No retained line contains a citation, so the wrapper carries no bibliography and no bibliography entry.  No retained line contains a figure environment, an image-inclusion command or drawing code, so no image is delivered and the package declares no asset dependency.  Editorial formatting only, in addition: an AMS \texttt{amsart} preamble that restates the article's own declarations and macros used by the retained lines, namely the environments \texttt{lemma} on separate counters, together with the standard packages those lines need.  The retained environments print in this document as Lemma 1 in order of appearance, whereas the article prints the same items with the numbering of its own full document.  The complete native source of the article as distributed in the arXiv e-print for this exact version is bundled unchanged as \texttt{sources/197206/source0.tex}.
    \begin{lemma}\label{lemma:technical_approx_limits}
        Let $u \in BV(\Omega; \R^M)$, $0 \in J_u$ with $\nu_{0} = e_n$, and $F \subset (0,1)$ be a set with $|F| = 1$. Then, there exists a sequence $\{\rho_m\} \subset F$ with $\rho_m \to 0^+$ such that 
        \[
            \H^{N-1}(J_u \cap \partial (-\rho_m/2, \rho_m/2)) = 0,
        \]
        and for the rescaled trace $u|_{\partial Q_{\rho_m}(x_0)}$ restricted to $(-\rho_m/2, \rho_m/2)\times \{\pm \rho_m\}$ we have 
        \[
            u(\rho_m \cdot, \pm \rho_m/2) \to u^\pm(x_0)
        \]
        in $L^1((-1/2, 1/2)^{N-1})$.
    \end{lemma}

    \begin{proof}
        Since $u^\pm(x_0)$ are one side approximate limits we have 
        \[
            \lim_{\rho \to 0^+} \int_{(0, 1)^N} |u(\rho x) - u^+(x_0)| \, dx = 0
        \]
        and
        \[
            \lim_{\rho \to 0^+} \int_{(-1, 0)^N} |u(\rho x) - u^-(x_0)| \, dx = 0.
        \]
        In particular, we have by Fubini's Theorem
        \[
            \lim_{\rho \to 0} \int_0^1 \int_{(0, 1)^{N}} |u(\rho x', \rho t) - u^+(x_0)| \, dx' \, dt= 0,
        \]
        where $u(\rho x', \rho t)$ agrees with the two-sided trace at $\partial (0,1)^N$ of $u$ that exists since 
        \begin{align}\label{eq:eligible_rho}
            \H^{N-1}(J_u \cap \partial (-\rho/2, \rho/2)) = 0
        \end{align}
        for all but countably many $\rho > 0$. In particular, for a set $E_\rho \subset (0,1)$ of full measure we obtain $t \in E_\rho$
        \[
            \lim_{\rho \to 0} \int_{(0, 1)^{N-1}} |u(\rho x', \rho t) - u^+(x_0)| \, dy' = 0,
        \]
        and, therefore, it follows that 
        \[
            \lim_{\rho \to 0} \int_{(0, 1)^{N-1}} |u(\rho t x', \rho t) - u^+(x_0)| \, dy' = \lim_{\rho \to 0} \frac1 {t^{N-1}} \int_{(0, t)^{N-1}} |u(\rho x', \rho t) - u^+(x_0)| \, dy' = 0.
        \]
        Now, choose a sequence $\{ \tilde \rho_m\} \subset (0,1)$ with $\tilde \rho_m \to 0$ for which \eqref{eq:eligible_rho} holds. Then, with an $t_m \in E_{\rho_m}$ such that $t_m\rho_m \in F$ holds, the sequence $\rho_m := t_m \tilde \rho_m$ satisfies the sought-after requirements.
    \end{proof}

\end{document}
